% !TeX root = main_long.tex
% RESTRUCTURED: one subsection per research question; incommensurable
% denominators kept apart; every figure traced to a pinned run or, for RQ4,
% to a digest-recorded artifact inspection.
\section{Results}
\label{sec:coverage}

\subsection{RQ1: does VCF Core preserve the meaning of the constructs in the VCF specification?}
\label{sec:evidence}

\paragraph{Version-specific declarations.}
The version registries hold 40, 40, 67, 78 and 122 explicit reserved declarations respectively for VCF 4.1--4.5; the 122 for VCF 4.5 comprise 51 INFO and 71 FORMAT definitions, three of them parameterized base-modification key families.
All \textbf{333} reserved Number/Type rows extracted from the specification texts (31, 31, 69, 79 and 123 per version) match those registries, with no missing or differing row.
Those 333 are the keys each source declares explicitly in \texttt{\#\#INFO=}/\texttt{\#\#FORMAT=} blocks or reserved-key tables, a different population from the registry entries: VCF 4.1 and 4.2 each define a further 29 reserved keys in prose bullets giving no Number and, for INFO, no Type, which are inventoried but not comparable, so the agreement figure is weighted towards VCF 4.3--4.5.
Matching declarations does not resolve inconsistencies between a declaration and its prose meaning: review found CNL/CNP declared \texttt{Number=G} alongside prose defining an ordering by copy number.

\paragraph{VCF specification source audit and review.}
The assessment proceeds in stages.
We first separated the VCF specificaiton latex{} source files into \textbf{344} source-section intervals, from which \textbf{1{,}134} passage-level assertions were authored.
An assertion becomes a \emph{requirement} only once it can be stated as a question with a varifiable answer, which yields \textbf{94} \emph{requirements}.
A \emph{requirement} becomes a \emph{case} when paired with one version and one \emph{fixture}, giving \textbf{210} total \emph{cases}.
Each \emph{case} includes a \emph{query} and an \emph{expected answer}.
Across the 210 \emph{cases}, we defined \textbf{79} distinct \emph{queries}, since one \emph{query} often serves several \emph{cases}.

Each of the 1{,}134 assertions is then classified by whether an executable test reaches it, using two authored fields: the \emph{cases} that test it, and a \emph{test-gap} note recording what those \emph{cases} leave untested.
An assertion with at least one \emph{case} and no recorded gap is \emph{targeted} (\textbf{149}), one with \emph{cases} but a recorded gap is \emph{partial} (\textbf{154}), and one with no \emph{cases} at all is \emph{untested} (\textbf{831}).
The distinction is visible in three real entries.
``Contig ID and supplied length belong to the declared reference sequence'' has one \emph{case} for R05 and no gap, so it is targeted.
``POS is 1-based and records retain distinct identity and source order even at equal coordinates'' has \emph{cases} for R07 and R08 but the authored gap ``coincident records need a targeted identity/order \emph{case}'', so its first half is tested and its second is not.
``The assembly URL identifies the external sequence source referenced by BKPTID values'' has no \emph{case} at all.
The links are checked rather than trusted: a cited \emph{case} must exist, share the assertion's version, name one of its \emph{requirements}, and carry a \emph{query} with an \emph{expected answer}.

These counts are not a coverage score, they only record whether a test exists, not whether it passed, and they are counted once per version, so an assertion stated identically in all five specifications is recorded five times.
The 1,134 assertions hold only 385 distinct statements, 86 of them verbatim in every version. 
Their purpose is to show where audit work has and has not reached; coverage is scored against the 94 \emph{requirements} instead.
Which assertions became \emph{requirements} was not arbitrary.
An \emph{expected answer} has to be defensible which is only practical where the specification is concrete enough to define an assertion and give an example.
Therefore, the investigation primarily concentrated on sections of the specification that provide examples: nine of the 38 \emph{fixtures} are byte-identical copies of published specification examples, and the remainder are minimal files written to instantiate a stated rule.
Assertions whose sources gave no defensible serialization and expected-answer basis remained unassessed.
Of the 94 \emph{requirements}, along the axis of "intent," 56 carry an interpretation specific enough to check the \emph{query} against a stated intent, 29 carry a generic intent, and 9 carry an explicit unassessed marker.

Deriving \emph{requirements} from the specification text also exposed defects in those examples, three of which were reported to the VCF specification maintainers via a github issue (\href{https://github.com/samtools/hts-specs/issues/869}{hts-specs\#869}, \href{https://github.com/samtools/hts-specs/issues/870}{hts-specs\#870}, and \href{https://github.com/samtools/hts-specs/issues/871}{hts-specs\#871}) and one issue already existed (\href{https://github.com/samtools/hts-specs/issues/868}{hts-specs\#868}). 
All issues remain open at the time of writing.
Each of these complicated authoring the affected \emph{queries} and \emph{expected answers} because the example that would ordinarily inform the \emph{oracle} was itself inconsistent with the assertion.
The affected \emph{requirements} were nevertheless scored rather than set aside.
Where it was not determinable because the passage is itself disputed upstream, no test was written and the \emph{requirement} was left unassessed, which counts against the score rather than being excluded from it.
The reported percentages are therefore depressed by unresolved specification inconsistencies rather than by a gap in the vocabulary.

\paragraph{Specification-derived \emph{requirements} (specification $\rightarrow$ vocabulary).}
The 210 \emph{cases} exercise 53 of the 94 registered \emph{requirements} in at least one applicable version, meaning that these 53 have defined tests. 
The remaining 41 have no registered \emph{case} in any version, meaning they do not yet have defined tests.
Coverage is reported per version, because \emph{requirements} apply per version (Table~\ref{tab:requirement-results}).
All 210 \emph{preservation} tests passed in each profile, as did all 210 \textit{expanded}-\emph{structure} tests.
In the \textit{condensed} profile 165 \emph{structure} tests passed and 45 failed, every failure concerns per-sample FORMAT access that the profile stores in encoded vectors rather than individual sample resources.
In VCF 4.5, 18 failed \textit{condensed}-\emph{structure} tests affected 14 \emph{requirements}, failed tests and affected \emph{requirements} are different counts.
Every passing \emph{case} rejected the empty-graph control and changed its answer under at least one referenced-predicate deletion.
The structural difference is therefore a sample-profile limitation in \textit{condensed} representations, since the corresponding \textit{expanded} tests passed it is not evidence that VCF Core cannot express those values structurally.

\begin{table}[htbp]
\centering
\caption{
  Demonstrated coverage of registered information \emph{requirements} by VCF version, \emph{sample profile}, and retrieval axis.
  $N$ is the total number of \emph{requirements} applicable to a given VCF version.
  ND (\emph{not demonstrated}) counts \emph{requirements} that were tested but did not pass, while
  U (\emph{unassessed}) counts \emph{requirements} for which no test was available; unassessed \emph{requirements} remain in $N$ and therefore contribute zero to the reported percentage. No \emph{requirement} received a partial result.
  \textbf{Each Demonstrated column is read on its own}: within a column, demonstrated $+$ ND $+$ U $= N$. For VCF 4.1 that is $32+0+37=69$ in the three agreeing columns and $26+6+37=69$ under \textit{condensed} \emph{structure}.
  Row totals are not meaningful, because the four columns report the same \emph{requirements} assessed four ways rather than four disjoint groups.
  $^{\dagger}$ND is nonzero only for the column immediately to its left, \emph{structure} under \textit{condensed}, and is zero in the other three combinations.
}
\label{tab:requirement-results}
\small
\begin{tabular*}{\linewidth}{@{\extracolsep{\fill}}lrrrrrrr@{}}
\toprule
 & & \multicolumn{4}{c}{Demonstrated} & & \\
\cmidrule(lr){3-6}
 & & \multicolumn{2}{c}{\emph{preservation}} & \multicolumn{2}{c}{\emph{structure}} & & \\
\cmidrule(lr){3-4} \cmidrule(lr){5-6}
VCF & $N$ & \textit{expanded} & \textit{condensed} & \textit{expanded} & \textit{condensed} & ND$^{\dagger}$ & U \\
\midrule
4.1 & 69 & 32 (46.4\%) & 32 (46.4\%) & 32 (46.4\%) & 26 (37.7\%) &  6 & 37 \\
4.2 & 72 & 34 (47.2\%) & 34 (47.2\%) & 34 (47.2\%) & 28 (38.9\%) &  6 & 38 \\
4.3 & 74 & 35 (47.3\%) & 35 (47.3\%) & 35 (47.3\%) & 29 (39.2\%) &  6 & 39 \\
4.4 & 83 & 38 (45.8\%) & 38 (45.8\%) & 38 (45.8\%) & 31 (37.3\%) &  7 & 45 \\
4.5 & 91 & 51 (56.0\%) & 51 (56.0\%) & 51 (56.0\%) & 37 (40.7\%) & 14 & 40 \\
\bottomrule
\end{tabular*}
\end{table}

\paragraph{How to read these numbers.}
Across versions, roughly half of the \emph{requirements} are demonstrated for the \textit{expanded} profile and about one third for the \textit{condensed} profile. 
Importantly, most remaining \emph{requirements} are \emph{untested rather than failed}: the not-demonstrated count is zero for every version, profile, and axis except \textit{condensed} structure. 
The appropriate interpretation is therefore that the reported percentages show what has been demonstrated, while the remainder is still unassessed.
For VCF 4.5, 40 \emph{requirements} remain unassessed. 
Of these, 37 await planned \emph{fixtures} and \emph{queries}, R45 needs further analysis because interpretation was unclear from the specification wording, and R67--R68 cover all reserved INFO and FORMAT keys and therefore cannot be satisfied by testing only selected \emph{fixtures}.

\paragraph{Curated inventory (vocabulary $\rightarrow$ specification).}
All \textbf{104} inventoried VCF 4.5 logical constructs were classified as both \emph{preserved} and \emph{structured}, with no partial or absent representations, in every specification area; \textbf{87} additionally had a rule checking a stated \emph{requirement}, enforcement being lower for file and header structure (17/20), lexical and encoding rules (6/7), FORMAT and genotype (16/21) and structural variation, repeats and gVCF (8/16)~\cite{vcfc210}.
The evidence behind those judgements was a \emph{query} for 7 constructs, a regression probe for 23, a shape or decoded check for 62 and a materialized \emph{fixture} for 12 --- evidence tiers behind curated judgements, not 104 independently executed retrieval tests.
Separately, all 17 \emph{fixtures} checked on the bytes pass, on their own axis.

All \textbf{104} inventoried VCF 4.5 logical constructs were classified as both \emph{preserved} and \emph{structured}, with no partial or absent representations across any specification area. 
Of these, \textbf{87} also had a validation rule enforcing a stated \emph{requirement}. 
Enforcement was less complete for file and header structure (17/20), lexical and encoding rules (6/7), FORMAT and genotype (16/21), and structural variation, repeats, and gVCF (8/16)~\cite{vcfc210}.
These classifications were supported by different forms of evidence: 7 constructs by \emph{queries}, 23 by regression probes, 62 by shapes or decoded checks, and 12 by materialized \emph{fixtures}. 
They should therefore be read as evidence supporting the curated inventory, rather than as 104 independently executed retrieval tests.


\subsection{RQ2: do RDF representations of VCF behave as documented?}
\label{sec:crossproducer}

\paragraph{The same questions, a second producer.}
RQ2 replicates RQ1 \emph{queries} and \emph{expected answer} assessments on \emph{witnesses} produced by VCF-RDFizer and compares the results to those produced by the VCF Core materializer. 
Out of 840 \emph{cases} tested, 693 pass for both producers and 45 fail for both; 53 \emph{requirements} are exercised, 38 of them agreeing on every check and 32 passing on every check.
The 45 both-fail checks are all \textit{condensed}-profile structure \emph{queries}: both producers answer the same way, and that way is ``no'' --- the documented storage trade-off, not a disagreement.
On 102 checks, spanning 15 \emph{requirements}, the materializer answers and the converter does not, and no check shows the inverse.
Because the paper's claim concerns the vocabulary rather than one program, we examined each of the 15 by running the failing \emph{query} against both graphs and checking which properties the converter emits (Table~\ref{tab:divergence}).


\begin{table}[htbp]
\caption{Why the converter fails a check the materializer passes. Every divergent \emph{requirement} is accounted for; the right-hand column is what would have to change to close it.}
\label{tab:divergence}
\centering\small
\begin{tabularx}{\linewidth}{@{}p{0.29\linewidth}cX@{}}
\toprule
Cause & Reqs. & What is missing \\
\midrule
A declared term the converter does not emit & 12 & The \emph{query} asks for a property the vocabulary defines and the converter never writes, so nothing matches. Twenty such properties, including gVCF reference blocks, breakend mates, allele indices, etc. \\
A decomposition that stops early & 1 & Tandem repeats: the summary row is emitted, the per-allele components are not. \\
A datatype the vocabulary leaves open & 1 & QUAL. The materializer writes \term{vcfc:VCFFloat}, the converter \term{xsd:decimal}; both satisfy the shape, and only a \emph{query} projecting \term{DATATYPE()} can tell them apart. \\
An encoding convention & 1 & An empty local-allele list, written \texttt{""} by one producer and \texttt{"."} by the other. \\
\bottomrule
\end{tabularx}
\end{table}

\textbf{None of the 15 divergences is the vocabulary being unable to express something.}
Across all 840 checks, none traces to a representation gap: every property involved is declared in the pinned vocabulary release, including the twenty the converter never writes.
What the divergences locate is unfinished or incorrect implementation --- 12 of the 15 are declared terms the converter does not yet emit --- so they bound what can be said about this converter, not about the model.

That makes the comparison a development instrument as much as a validation one, because it sees what a producer's own tests cannot.
A test suite compares a producer against itself, so it cannot detect a value that is well-formed, plausible, and attached to the wrong thing.
An earlier independent RDF converter comparison found exactly that: the converter bound \term{Number=LA} and \term{Number=LR} items to allele indices positionally, as if the list indexed the record's ALT column, when it indexes the sample's local allele subset.
Under \texttt{LAA=2,4} a depth list of \texttt{20,30,10} belongs to global alleles 0, 2 and 4; the converter attached it to 0, 1 and 2 --- wrong values rather than missing ones, with nothing in the graph to signal it.
That defect and eight further divergences were repaired in response~\cite{vcfrConverter2026}, and the pinned release additionally emits the ID, FILTER, field-index and per-allele phase-indicator accessors it had never written, and recognises the thirty base-modification aliases it had ignored.
All local-allele \emph{requirements} now pass for both producers on every \textit{expanded}-profile check; one \textit{condensed}-profile divergence remains, R27, described next.

% TODO: maybe remove this paragraph
\textbf{None of the 15 divergences indicates a limitation of the vocabulary itself.}
Across all 840 checks, no discrepancies trace to a representation gap. 
The differences identified above instead arise from the implementation of the independent converter VCF-RDFizer: for 12 of the 15 affected \emph{requirements}, the required terms exist in the vocabulary but were not yet emitted by the converter. 
These results therefore limit what can be claimed about the converter, rather than what the vocabulary can represent.
That makes the comparison performed a development instrument as well as a validation one, because it reveals what a producer's own tests cannot.


\paragraph{Semantic round-tripping.}
Across all 38 \emph{fixtures}, translated into \emph{wittnesses} by both conversion mechanisms, reconstructing a VCF file from structured properties alone recovers all \textbf{109 records} on the seven fixed columns and all \textbf{151 sample genotypes}; the companion check recovers all \textbf{431 header lines}, compared by source line number on key and verbatim value, and \textbf{400 of 401} INFO entries and non-GT FORMAT subfields.
Information therefore survives conversion into the structured layer and can be read back out of it without touching a preserved source string.
The one exception is the \texttt{LAA} field of a record whose local-allele list is empty.
In this case, the materializer emits the field resource, its \term{declaredBy} link and \term{fieldValue ""}, while the converter emits the resource and the link but no value triple, so the reconstructed cell reads \texttt{0/0:?:30:0}.
An empty \texttt{LAA} is meaningful rather than missing, it says this sample's local alleles are the reference only (genotype \texttt{0/0}), so what is lost is the distinction between \emph{present and empty} and \emph{present without a value}.
This is \emph{requirement} R27, reached here by a different route than the independent materialization took: two checks landing on one defect from opposite directions.
It is not a representation gap, since the vocabulary expresses the empty value and the materializer writes it, and it does not propagate, since \texttt{LAD}'s value item still carries its \term{forAllele} link on the same record.

\paragraph{Where the two \emph{sample profiles} (\textit{condensed} and \textit{expanded}) agree.}
\label{sec:profilesresult}
The \emph{sample profiles} return the same outcome on \textbf{352 of the 420} \emph{case}-and-axis pairs.
All 68 disagreements fall within 15 \emph{requirements} and share one property: every one reads data held \emph{inside} a sample, either a FORMAT value or a genotype.
Questions about the file, its header, a record, or an allele agree on every check in both profiles.
This discrepancy is infact the \textit{condensed} profile working as designed. 
Becasue the \textit{condensed} profile stores sample values as vectors rather than one resource per sample and field, a \emph{query} written to walk per-sample resources finds nothing, and recovering the value requires a decoding step the included \emph{queries} refuse.
The supported claim is therefore narrower than ``the profiles are equivalent'': they agree above the sample level and diverge by design below it.

% TODO: I think this paragraph can go as well
\paragraph{VCF Core version-level checks.}
Converting the worked example emits 242 \textit{expanded} and 168 \textit{condensed} triples, drawing on 93 and 80 distinct VCF Core terms respectively, and every one of those terms is declared in VCF Core 2.1.2 as resolved from the git tag.
Separately, the vocabulary's own recorded validation run accepts 21 \emph{fixtures} and includes a negative control, alongside a regression suite of 101 tests.
Because the converter pins no vocabulary version, the relationship between the two is established by the recorded behaviour rather than by a declaration, and it is evidence for the terms these graphs exercise rather than for all 467 declared terms or every historical VCF version.

\subsection{RQ3: does the representation support an integration task?}
\label{sec:integration}

Asking which flagged alterations have supporting read depth of at least 20 returns two rows, identical to the answers authored in advance: \texttt{chr1:1~G>C} at depth 30 and \texttt{chr1:2~A>T} at depth 25, each traced to its file, record, sample and the FORMAT declaration giving \texttt{LAD} its cardinality, with every row naming the \emph{fixture} actually converted.
The three flagged alterations absent from the answer are each confirmed absent individually and for a stated reason: \texttt{chr1:1~G>A} and \texttt{chr1:2~A>C} because ALT 1 is not in the sample's \texttt{LAA} subset at those positions, and \texttt{chr1:3~C>G} because ALT 1 is not in \texttt{LAA=3} there and its positional value would fall below the threshold in any case.

This example importantly highlights that a "shallower" representation (i.e. one that cannot semantically represent file, record, sample, and FORMAT fields) gives a different answer than the VCF Core representation.
If \texttt{Number=LR} values are interpreted only by position, the $n$th \texttt{LAD} value is matched to the $n$th ALT allele in the record rather than to the $n$th allele in the sample's \texttt{LAA} subset. 
Under this interpretation, the \emph{query} still returns two rows, but for the wrong alterations: \texttt{chr1:1~G>A} at depth 30 and \texttt{chr1:2~A>C} at depth 25.
The result therefore appears structurally plausible: it has the expected number of rows and retains provenance, but the values are attached to the wrong alleles. 
The expected-answer set records this positional interpretation separately, using only the \emph{fixture} and specification rule defined before \emph{query} execution. 
Two of the three exclusions change only when \texttt{Number=LR} is resolved through the sample-specific \texttt{LAA} mapping. 
The difference therefore comes from how the VCF semantics are represented, not simply from whether the underlying values are present.
This comparison illustrates the consequence of the specification rule rather than measuring an error in a particular converter. 
The example is intentionally limited: the input is synthetic, contains four sites and one sample, and uses invented annotations with no biological or clinical meaning. 
It demonstrates the importance of local-allele-aware modelling for this \emph{fixture}, not correctness across all possible VCF data.


\subsection{RQ4: how does VCF Core relate to other VCF semantic models?}
\label{sec:models}

Table~\ref{tab:models} records the nine verdicts per artifact; three patterns account for most of it.

\emph{HERO-Genomics and GVO keep INFO as one literal.}
Both \term{hero:vcfInfo} and \term{gvo:info} are datatype properties carrying the whole INFO column as a single string, and both align to \term{vcfc:infoRaw}, so the information survives but every INFO-derived question requires re-parsing that string \emph{and} supplying the cardinality rule from outside the graph --- which is why F7 is C rather than E for both.
\emph{GFVO made the opposite choice}, modelling particular field \emph{meanings} as first-class concepts (\term{AlleleCount}, \term{Coverage}, \term{MappingQuality}), which answers those fields well and leaves the declaration mechanism unaddressed, so F2 is S.

\emph{VCF2RDF is the instructive case, and it corrected an earlier verdict of ours.}
Reading the paper alone suggested a term-free isomorphic mapping; the published artifact shows otherwise, with 58 dereferenceable term pages including \term{Number}, \term{Type} and \term{Description}.
But with no allele-index term a consumer must still split the value list and count ALT alleles by re-reading a literal: having the declaration is necessary and, on its own, insufficient.

\begin{table}[htbp]
\caption{
  Can each model's published artifact answer the question? \textbf{E}: explicitly represented. \textbf{C}: representable with additional conventions the model does not itself establish. \textbf{S}: outside the model's stated scope --- a difference of purpose, not a defect. No cell remained unsettled.
}
\label{tab:models}
\centering\small
\begin{tabularx}{\linewidth}{@{}Xccccc@{}}
\toprule
Can a consumer\ldots & VCF Core & GFVO & HERO & GVO & VCF2RDF \\
\midrule
F1\quad trace a record to its file & E & E & E & S & C \\
F2\quad read a declaration's ID, Number, Type as data & E & S & S & S & E \\
F3\quad read the declared VCF version & E & S & S & S & C \\
F4\quad address ALT allele $n$ by index & E & C & C & S & C \\
F5\quad separate ploidy, allele calls and phasing & E & C & C & S & C \\
F6\quad get a (sample, field) value unsplit & E & C & C & S & C \\
F7\quad bind a \texttt{Number=A/R/G} item to its allele & E & C & C & S & C \\
F8\quad distinguish explicit \texttt{.} from ``not stated'' & E & S & C & C & C \\
F9\quad tell which sample encoding was used & E & S & S & S & S \\
\bottomrule
\end{tabularx}
\end{table}

\emph{GVO's column is almost entirely S} because the retrieved ontology is primarily organized around classes for types of genomic variation. 
It declares 48 classes (47 of which represent variation types beneath the root class \term{gvo:Variation}) and 12 datatype properties, but no object properties.
As a result, the ontology does not provide relationships for linking variation records to resources such as alleles or samples.
This should not be read as a deficiency. 
GVO is designed to classify genomic variation, so the S entries reflect a difference in modelling purpose rather than missing functionality relative to its intended scope.
The only C concerns missing values. 
A literal-valued property can contain \texttt{.}, but the ontology itself does not define how this should be distinguished from the complete absence of a corresponding triple.


\paragraph{What this supports, and what it does not.}
\emph{Supported}: VCF Core is distinguished by its representation of the \emph{interpretive context} carried by a VCF file, not by broader modelling of variation itself --- GVO's 47 variation classes are more extensive in that respect. 
This includes the declaration that determines a field's cardinality, the VCF version that determines which rule applies, the allele index needed to resolve genotype values, and the sample column that defines the scope of a value.

In the other inspected ontologies this context is handled in one of three ways, each visible in the matrix.
HERO and GVO keep the INFO column as a single literal --- \term{hero:vcfInfo} and \term{gvo:info} both align to \term{vcfc:infoRaw} --- so the text survives, but every INFO-derived question requires parsing that string again and supplying the cardinality rule from outside the graph.
GFVO makes the opposite choice, promoting particular field \emph{meanings} to first-class concepts such as \term{AlleleCount} and \term{MappingQuality}, which serves those fields well while leaving the declaration mechanism itself unaddressed.
VCF2RDF does publish the declaration, with terms for \term{Number}, \term{Type} and \term{Description}, but offers no allele index, so a consumer must still split the value list and count ALT alleles by re-reading a literal.

\emph{Not supported}: these results do not show that the other models are incapable of representing such information. 
A label of ``requires additional conventions'' refers only to what is explicitly established by the inspected ontology artifacts, not to what its community could extend or implement.

Reuse and cross-model alignment are therefore reported separately. 
VCF Core directly reuses terms from PROV, DCAT, SO, and FALDO, and provides 123 curated SSSOM mappings~\cite{matentzoglu2022sssom} to variation models and reference ontologies with different levels of confidence, see \href{vcf-core-vocabulary/mappings}{https://github.com/ecrum19/vcf-core-vocabulary/tree/main/mappings}. 
These mappings form a separate alignment artifact and should not be interpreted using the same denominator as the comparison matrix above~\cite{vcfc210}.

